\documentclass[12pt]{scrartcl}

%Matthis Vorlage
\usepackage{../mathe_definition}
\usepackage[shortlabels]{enumitem}

%Ausgelagerter Code fürs schöne seitenlayout
\date{\today}
\newcommand{\Blattnummer}{%
	3%
}

\title{{\itshape Stochastische Prozesse}}
\module{Blatt 1}
\author{Philip Wunderlich \and Jonathan Langer}
%\matriculation{3737927\and 3750030}

\fancypagestyle{netter_header}{
	\renewcommand{\and}{\\}
	\fancyhead[L]{\MyAuthor}
	\fancyhead[C]{\MyTitle\\ \MyModule}
	\fancyhead[R]{\MyDate\\Seite \thepage{} von \pageref{LastPage}}
	\fancyfoot[C]{}
}
\pagestyle{netter_header}

%Abkürzungen
\newcommand{\einruecken}{\begin{addmargin}[25pt]{0pt}} 
	%LaTeX mag keine Umlaute :(
	\newcommand{\ausruecken}{\end{addmargin}}
\newcommand{\R}{\ensuremath{\mathbb{R}}}
\newcommand{\Om}{\ensuremath{\Omega}}
\newcommand{\om}{\ensuremath{\omega}}
\newcommand{\A}{\ensuremath{\mathcal{A}}}
\newcommand{\B}{\ensuremath{\mathcal{B}}}
\newcommand{\C}{\ensuremath{\mathcal{C}}}
\newcommand{\WraumAnn}{\ensuremath{(\Om, \A, P)} ein Wahrscheinlichkeitsraum}
\newcommand{\MraumAnn}{\ensuremath{(E,\B)} ein Messraum}
\newcommand{\Xt}{\ensuremath{(X_t)_{t\in T}}}
\newcommand{\Yt}{\ensuremath{(Y_t)_{t\in T}}}
\newcommand{\Xw}{\ensuremath{X_t(\omega)}}
\newcommand{\Yw}{\ensuremath{Y_t(\omega)}}
\newcommand{\Xq}{\ensuremath{X_q(\omega)}}
\newcommand{\Yq}{\ensuremath{Y_q(\omega)}}
\newcommand{\E}{\ensuremath{\mathbb{E}}}
\newcommand{\1}{\ensuremath{\textbf{1}}}
\newcommand{\LwieLp}{\ensuremath{\mathcal{L}}}
% \L gibts anscheinend schon
\newcommand{\scriptL}{\ensuremath{\mathcal{L}}}
\newcommand{\scriptP}{\ensuremath{\mathcal{P}}}
\newcommand*\diff{\mathop{}\!\mathrm{d}}


\begin{document}
	
\subsection*{Aufgabe 10.}

Sei \WraumAnn{}, $\F=(\F_n)_{n\ge 0}$ eine Filtration und $X=(X_n)_{n\ge 0}$ ein an \F{} adaptierter Stochastischer Prozess mit $X_n\in\scriptL_1(P)$ für alle $n\ge 0$.

\subsubsection*{Beweis}
\einruecken

trivial

\ausruecken
\subsection*{Aufgabe 11.}

Sei $(M_n)_{n\ge 0}$ ein $(\F_n)_{n\ge 0}$-Martingal mit $M_0=0$ für das es Konstanten $c_n\ge 0$ gibt, so dass $|M_n-M_{n-1}|\le c_n$ fast sicher für alle $n\ge 1$ gilt. Zeige
$$P(|M_n|\ge\lambda)\le 2\exp \left(-\frac{\lambda^2}{2\sum_{i=1}^n c_i^2}\right),\quad \lambda\ge 0, n\ge 1.$$
Anleitung:
\begin{enumerate}[$i)$]
\item Betrachte $D_n:=M_n-M_{n-1}$ für $n\ge 1$. Benutze die Konvexität der Exponentialfunktion und die Taylorreihenentwicklung von $\cosh(\lambda c_n)$, um für $t>0$ zu zeigen:
$$\E(e^{tD_n}|\F_{n-1})\le\exp(t^2 c_n^2/2).$$
\item \label{11.ii} Verwende die Ungleichung aus $i)$, um für $t>0$ induktiv zu zeigen:
$$\E\left(\exp\left(t\sum_{i=1}^n D_i\right)\right)\le\exp\left(t^2\sum_{i=1}^n c_n^2/2\right).$$
\item Verwende die Markov-Ungleichung und die Abschätzung aus \ref{11.ii} für $t>0$ beliebig. Optimiere anschließend bezüglich $t$.
\end{enumerate}

\subsubsection*{Beweis}
\einruecken

trivial

\ausruecken
\subsection*{Aufgabe 12}

Seien $(D_n)_{n\ge 1}$ i.i.d. mit $P(D_1=-1)=P(D_1=1)=1/2$. Wir interpretieren jedes $D_i$ als Ergebnis einer fairen Wette, bei der im Misserfolgsfall der Einsatz in der nächsten Runde verdoppelt wird. Im Erfolgsfall wird dagegen das Spiel beendet. Dies beschreibt die folgende Verdopplungsstrategie
$$H_n:=2^{n-1}\1_{\{D_1=\dots=D_{n-1}=-1\}}, \quad n\ge 1$$
für den Einsatz in der $n$-ten Runde. Für diese Spielstrategie ist der Gesamtgewinn nach $n$ Runden durch
$$G_n=\sum_{i=1}^n H_iD_i$$
gegeben.

\begin{enumerate}[$i)$]
\item Finde ein Martingal $X=(X_n)_{n\ge 0}$ für das gilt $G=H.X$.
\item Zeige $\E G_n=0$ für alle $n\ge 1$.
\item Zeige $\lim_{n\rightarrow\infty}G_n=1$ fast sicher.
\end{enumerate}

\subsubsection*{Beweis}
\einruecken

trivial

\ausruecken

\end{document}